\chapter{Organizzazione del laboratorio}
\label{cap:organizzazione}

\section{Percorso essenziale in quattro laboratori}

\begin{center}\small
\begin{tabular}{lp{5.2cm}p{6.6cm}}
\toprule
 & Contenuto & Obiettivi di apprendimento \\
\midrule
Lab 1 & Produzione e scorte (cap.~\ref{cap:produzione}) &
formulare un LP multiperiodale; leggere duali, slack e range; verificare un prezzo
ombra per perturbazione \\
Lab 2 & Markowitz (cap.~\ref{cap:markowitz}) &
costruire un QP convesso; tracciare una frontiera; discutere la fragilità delle stime \\
Lab 3 & Pricing \emph{oppure} budget (capp.~\ref{cap:pricing}--\ref{cap:budget}) &
modellare funzioni non lineari; studiare concavità; verificare le KKT numericamente \\
Lab 4 & Progetto a scelta &
supply chain, ricarica EV, localizzazione, code, Newsvendor, CVaR o SVM; presentazione
manageriale dei risultati \\
\bottomrule
\end{tabular}
\end{center}

I capitoli~\ref{cap:newsvendor}--\ref{cap:svm} possono sostituire i laboratori 3--4 in
un'edizione orientata all'incertezza e al machine learning, oppure diventare
approfondimenti per progetti individuali.

\section{Struttura della consegna}

Ogni gruppo consegna un notebook o un breve report (max 8 pagine) con questa
struttura obbligatoria, che ricalca la scaletta dei capitoli:

\begin{enumerate}
  \item \textbf{Problema e ipotesi} --- il contesto in una pagina, le semplificazioni
        dichiarate esplicitamente;
  \item \textbf{Modello} --- insiemi, parametri, variabili, vincoli, obiettivo, con la
        spiegazione di ogni vincolo (non solo la formula);
  \item \textbf{Dati} --- origine, unità di misura, eventuale generazione;
  \item \textbf{Risultati} --- valore ottimo, decisioni principali, vincoli attivi;
  \item \textbf{Sensitività} --- protocollo del capitolo~\ref{cap:richiami},
        sezione~\ref{sec:protocollo}: almeno prezzi ombra verificati, una
        one-at-a-time e una frontiera di trade-off;
  \item \textbf{Raccomandazione manageriale} --- massimo dieci righe, senza formule:
        cosa fare, quanto vale, quanto è robusto.
\end{enumerate}

\section{Criteri di valutazione}

\begin{center}\small
\begin{tabular}{lc p{7.6cm}}
\toprule
Dimensione & Peso & Che cosa si valuta \\
\midrule
Correttezza della formulazione & 30\% & vincoli giusti e giustificati, unità coerenti,
convessità discussa \\
Implementazione e verifica & 25\% & codice riproducibile; stato del solver
controllato; almeno una verifica indipendente (esempio a mano, perturbazione,
caso limite) \\
Analisi di sensitività & 25\% & protocollo completo; duali interpretati con i range
di validità; trade-off tracciati \\
Interpretazione e comunicazione & 20\% & raccomandazione chiara; figure leggibili;
onestà sui limiti del modello \\
\bottomrule
\end{tabular}
\end{center}

\section{Domande tipiche di discussione}

Le domande con cui difendere la consegna --- e che conviene porsi \emph{prima}:

\begin{itemize}
  \item Quale risorsa conviene aumentare per prima, e quanto si può pagare per essa?
  \item Qual è il costo di una promessa di servizio più ambiziosa?
  \item La soluzione resta credibile se i dati cambiano del 5\%? Quale parametro è
        il più fragile?
  \item Quale compromesso consigliereste a un decisore, e perché proprio quel punto
        della frontiera?
  \item Che cosa NON dice il modello? (ipotesi, orizzonte, variabili escluse)
\end{itemize}

\section{Errori ricorrenti da evitare}

\begin{attenzione}[Checklist degli errori visti più spesso]
\begin{enumerate}
  \item Leggere \texttt{.X} o \texttt{.Pi} senza controllare \texttt{m.Status}.
  \item Dimenticare \texttt{lb=-GRB.INFINITY} sulle variabili libere
        ($b$ della SVM, $\eta$ del CVaR).
  \item Interpretare un prezzo ombra fuori dal suo intervallo di validità
        (\texttt{SARHSLow/Up}).
  \item Confondere il segno dei duali nei problemi di minimo (vincolo $\le$:
        $\texttt{Pi} \le 0$).
  \item Aggiornare il RHS di un vincolo che contiene costanti a sinistra
        (capitolo~\ref{cap:ricaricaev}).
  \item Ottimizzare un solo obiettivo quando il problema ne ha due (minimax puro
        senza costo).
  \item Scegliere gli iperparametri ($C$, $\gamma$, $\lambda$) guardando il test set
        o gli scenari usati per decidere.
  \item Riportare sei cifre decimali da stime che ballano alla seconda.
\end{enumerate}
\end{attenzione}

\section{Materiali del corso}

\begin{center}\small
\begin{tabular}{lp{8.6cm}}
\toprule
\texttt{dispensa/} & questa dispensa (\texttt{main.tex}; le figure leggono i CSV in
\texttt{figure/dat/}) \\
\texttt{python/} & uno script per capitolo; \texttt{esegui\_tutti.py} rigenera
dati e figure \\
\texttt{dati/} & tutti i dati dei casi di studio in CSV \\
\texttt{GUIDA\_GUROBI.md} & guida autonoma al solver \\
\texttt{soluzioni/} & soluzioni degli esercizi (per i docenti) \\
\texttt{slides/} & slide dei quattro laboratori (italiano e inglese) \\
\bottomrule
\end{tabular}
\end{center}

Riproducibilità: \texttt{python3 -m pip install gurobipy matplotlib pandas scipy},
poi \texttt{python3 python/esegui\_tutti.py} e \texttt{latexmk -pdf dispensa/main.tex}.
